\section{Summable directional pinching}
\label{pi:section}

We prove the estimate which relates a difference of coordinate values to
the Bochner measure.  The summable quantity in this section is an energy;
its square root will occur in individual estimates, but will never be
summed without an exponentially decreasing weight.

\subsection{Conventions and domains}

We use the sequence of rescaled spaces and harmonic columns constructed
in Section~\ref{co:section}.  We suppress the index of that sequence.
Thus the lower Ricci bound is $-\kappa$, with $\kappa\to0$, and $o(1)$
is uniform over the centers and dyadic scales under consideration.
The measure is always the Hausdorff measure in the current length units.
We write $r_i=2^{-i}$, $i\ge0$, and use the matrices $R_i$ and the
matrix-valued positive measure $M(u)$ from the preceding section.

Fix numbers $L,C_{\mathrm{sep}}>1$.  The final estimate will involve the mass on
$B_{Lr_i}(y)$ and centers satisfying
\begin{equation}
 x,z,y\in B_4(p),\qquad
 d(x,y),d(z,y)\le C_{\mathrm{sep}}r_i.
 \label{pi:center-configuration}
\end{equation}
There is no restriction in replacing an initially prescribed number by
a larger number greater than one.  Set
\begin{equation}
 R=16(L+C_{\mathrm{sep}}+2),\qquad S=32R.
 \label{pi:radii}
\end{equation}
All the balls used below are contained in $B_{16S+4}(p)$ before the
additional dyadic rescaling.  We choose this fixed enlargement first
and then take a sufficiently late member of the sequence.  In
particular, the harmonic functions are defined there and all the
conclusions of Lemma~\ref{co:harmonic-normalization} apply on its
fixed enlargements.  Poisson solutions centered at $a$ will be defined
on $B_{8Sr_i}(a)$.  Their estimates will be used on $B_{2Sr_i}(a)$;
the affine fits will be made on $B_{Sr_i}(a)$; radial invariance will
be used only on $B_{Sr_i/2}(a)$.  In the two-center argument, its
projection ball is $B_{Rr_i}(z)$ and the transport calculation uses
$B_{3Rr_i}(z)$.  These choices leave strict margins between every
testing ball and its domain of definition.

For a measurable ball $B$ of positive measure, put
\[
 \|f\|_{2,B}^{\mathrm{av}}
 =\left(\fint_B |f|^2\dd\m\right)^{1/2}.
\]
Unadorned $L^2$ norms below use integrals, rather than averages.
After a dyadic rescaling, the volumes of all fixed-radius balls in
use are bounded above and below by positive constants.  Thus these
two norm conventions differ by uniformly bounded factors at each of
the fixed radii.

We use two consequences of the preceding section.  First, the
Dirichlet inequalities on the rescaled balls in use include, for
some $d>1$ and a constant depending on the fixed radius,
\begin{equation}
 \|\psi\|_{L^{2d}}^2\le C\int |\nabla\psi|^2\dd\m,
 \qquad \psi\in W^{1,2}_0(B).
 \label{pi:dirichlet-sobolev}
\end{equation}
as provided by Lemma~\ref{co:local-calculus}, specifically
\eqref{co:dirichlet-sobolev}.
The corresponding Poincar\'e inequality is included.  In dimension
two one may choose any one of the available exponents strictly
greater than two.  The domains are proper subsets of larger balls
whose complements have a fixed positive volume, so the Dirichlet
inequality is coercive.  Second, if
\[
 \sigma_i^a(b)=|b|_{R_i(a)},\qquad b\in\R^n,
\]
then, after fixing for example $a_0=1/16$ in the growth estimate,
\begin{equation}
 \frac{\sigma_j^a(b)}{\sigma_i^a(b)}
 \le C2^{a_0|i-j|}\quad(b\ne0),\qquad
 \sup_{B_{2Sr_i}(a)}|\nabla(b\cdot u)|
 \le C\sigma_i^a(b).
 \label{pi:coordinate-input}
\end{equation}
Each constant in this section may depend on $n,L,C_{\mathrm{sep}}$, and consequently
on $R,S$, but not on $i$, the centers, or the late member of the sequence.

\subsection{Positive sources and Poisson pinching}

Poisson replacements of squared distance follow the almost-rigidity
method of Cheeger and Colding \cite{CC96}; see also the quantitative
regularization in \cite[Section~6.1]{CJN21}. We retain the full
positive Bochner measure and prove the uniform estimates needed at
all testing centers.

The following elementary Dirichlet-form observation is useful because
an application of Cauchy--Schwarz to the energy estimate alone would
not give the required dependence on the mass.

\begin{lemma}[A linear gradient bound for a positive source]
\label{pi:positive-source}
Let $\Omega$ be one of the bounded rescaled domains satisfying
\eqref{pi:dirichlet-sobolev}, and suppose that $h\in W^{1,2}_0(\Omega)$
is nonnegative and bounded, with
\[
 \int\langle\nabla h,\nabla\psi\rangle\dd\m
 =\int\psi\dd\mu
\]
for compactly supported Lipschitz tests, where $\mu$ is a finite
positive measure.  If $\mu(\Omega)\le P$, then
\begin{equation}
 \int_\Omega|\nabla h|\dd\m\le CP,
 \qquad
 \int_\Omega|\nabla h|^2\dd\m\le\|h\|_\infty P.
 \label{pi:positive-source-bounds}
\end{equation}
The constant in the first inequality depends only on the Sobolev
constant, its exponent, and an upper bound for $\m(\Omega)$.
\end{lemma}

\begin{proof}
We first explain the interpretation of tests against $\mu$.  Since
$\mu=-\Delta h$, its action extends continuously to
$W^{1,2}_0(\Omega)$.  This positive functional does not charge a set
of zero Sobolev capacity: for a compact such set, test it with
nonnegative functions which are at least one near that set and whose
$W^{1,2}_0$ norms tend to zero, and use continuity of the functional.
Regularity of the measure gives the same conclusion for a Borel set
of zero capacity.  Consequently bounded Sobolev functions are
integrated using their quasi-continuous representatives.  Approximation
in $W^{1,2}_0$, followed by bounded truncation and a subsequence
converging outside a set of zero capacity, shows that the functional
and the measure integral agree on these tests.  In particular all
bounded Lipschitz compositions of $h$ which vanish at zero are
admissible.  This justifies the truncations used next without making
a choice of arbitrary $\m$-almost-everywhere representatives on the
support of $\mu$.

If $P=0$, testing with $h$ proves $h=0$.  Suppose $P>0$ and put
$H=h/P$.  For $t>0$, the test $\min\{H,t\}$ gives
\[
 \int_{\{H<t\}}|\nabla H|^2\dd\m\le t.
\]
Applying \eqref{pi:dirichlet-sobolev} to the same truncation yields
\[
 t^2\m(H\ge t)^{1/d}\le Ct,
 \qquad \m(H\ge t)\le Ct^{-d}.
\]
Choose $0<\alpha<d-1$.  The layer-cake formula, using the volume
bound for $t\le1$ and the last estimate for $t\ge1$, implies
\begin{equation}
 \int_\Omega(1+H)^{1+\alpha}\dd\m\le C.
 \label{pi:source-moment}
\end{equation}
The bounded test $1-(1+H)^{-\alpha}$ gives
\[
 \alpha\int_\Omega
 \frac{|\nabla H|^2}{(1+H)^{1+\alpha}}\dd\m\le1.
\]
Combining this inequality with \eqref{pi:source-moment} by weighted
Cauchy--Schwarz proves $\int|\nabla H|\le C$, which is the first
assertion.  Finally $h$ itself is an admissible bounded test, so
$\int|\nabla h|^2=\int h\dd\mu\le\|h\|_\infty P$.
\end{proof}

For a function with constant Laplacian and a scalar $b$, write
\begin{equation}
 \mathcal B_\kappa(w,b)
 =M(w)-(2b\Delta w-nb^2)\m.
 \label{pi:shifted-measure}
\end{equation}
By Lemma~\ref{co:local-calculus}, this is a nonnegative measure and
\begin{equation}
 \mathcal B_\kappa(w,b)
 \ge |\Hess w-bg|^2\m.
 \label{pi:shifted-lower-bound}
\end{equation}
It is quadratic in the pair $(w,b)$.  Positivity, applied to the sum
and difference of two such pairs, gives the measure inequality
\begin{equation}
 \mathcal B_\kappa(w_1-w_2,b_1-b_2)
 \le2\mathcal B_\kappa(w_1,b_1)
    +2\mathcal B_\kappa(w_2,b_2).
 \label{pi:shifted-triangle}
\end{equation}
These statements use the noncollapsed dimension and trace identity
and the improved Bochner inequality; see \cite{DPG18,Han18,Gigli18}.

\begin{lemma}[Poisson replacements with summable defects]
\label{pi:poisson-pinching}
There are nonnegative numbers $P_i(a)$, defined for every
$a\in B_4(p)$ and $i\ge0$, with
\begin{equation}
 \sup_{a\in B_4(p)}\sum_{i=0}^{\infty}P_i(a)=o(1),
 \label{pi:poisson-sum}
\end{equation}
and functions $f_i^a$ on $B_{8Sr_i}(a)$ such that, with
$q^a=d_a^2/2$,
\begin{equation}
 \Delta f_i^a=a_i:=n+C_0\kappa r_i^2,
 \qquad f_i^a-q^a\in W^{1,2}_0(B_{8Sr_i}(a)).
 \label{pi:poisson-definition}
\end{equation}
Here $C_0$ depends only on $n,S$ and is the same for every center
and every scale.  The following bounds hold:
\begin{align}
 0\le q^a-f_i^a&\le Cr_i^2,
 &\sup_{B_{4Sr_i}(a)}|\nabla f_i^a|&\le Cr_i,
 \label{pi:poisson-boundedness}\\
 \int_{B_{8Sr_i}(a)}|\nabla(f_i^a-q^a)|^2\dd\m
 &\le Cr_i^{n+2}P_i(a),
 &\int_{B_{8Sr_i}(a)}|\nabla(f_i^a-q^a)|\dd\m
 &\le Cr_i^{n+1}P_i(a),
 \label{pi:poisson-errors}\\
 \mathcal B_\kappa(f_i^a,1)(B_{2Sr_i}(a))
 &\le Cr_i^nP_i(a).
 \label{pi:poisson-shifted-mass}
\end{align}
\end{lemma}

\begin{proof}
Choose $C_0$ sufficiently large that Laplacian comparison gives
\[
 \Delta q^a\le(n+C_0\kappa r_i^2)\m
 \quad\hbox{on }B_{16Sr_i}(a).
\]
This is possible because $d_a\le16Sr_i$ there.  Take the sequence
late enough that $C_0\kappa\le1$.  Define the positive measure
\[
 \mu_i^a=a_i\m-\Delta q^a
 \quad\hbox{on }B_{16Sr_i}(a),
 \qquad
 P_i(a)=r_i^{-n}\mu_i^a(B_{8Sr_i}(a))+\kappa r_i^2.
\]
Here the measure comparison for the squared distance is on the
whole ball, including its center; see
\cite[Corollary~5.15, equation~(5.29)]{Gigli15}.
The source restricted to $B_{8Sr_i}(a)$ is a continuous functional
on $W^{1,2}_0$ since it is the difference of a bounded density and
the distributional Laplacian of the Sobolev function $q^a$.
The Dirichlet variational problem therefore gives
$h_i^a\in W^{1,2}_0$ with $-\Delta h_i^a=\mu_i^a$.
Testing its negative part gives $h_i^a\ge0$, and
$f_i^a=q^a-h_i^a$ satisfies \eqref{pi:poisson-definition}.

Here is the mass estimate, including its normalization.  Set
$V_a(t)=\m(B_t(a))$ and $A_a(t)=t^{-n}V_a(t)$.  Radial testing of
the measure-valued Laplacian of $q^a$, with $|\nabla d_a|=1$
almost everywhere off $a$, gives
\[
 \Delta q^a(B_t(a))=tV_a'(t)
 \quad\hbox{for almost every }t.
\]
Equivalently, this follows by approximating radial indicators in
$-\int t\psi'(t)V_a'(t)\dd t$; the boundary term at zero vanishes
by the local volume upper bound.  Thus
\[
 \mu_i^a(B_t(a))=(a_i-n)V_a(t)-t^{n+1}A_a'(t)
\]
at those radii.  Positivity and integration for
$8Sr_i<t<16Sr_i$ imply
\begin{align}
 r_i^{-n}\mu_i^a(B_{8Sr_i}(a))
 &\le C\int_{8Sr_i}^{16Sr_i}
       \frac{\mu_i^a(B_t(a))}{t^{n+1}}\dd t\notag\\
 &\le C\{A_a(8Sr_i)-A_a(16Sr_i)+\kappa r_i^2\}.
 \label{pi:volume-drop}
\end{align}
The last step uses $V_a(t)\le Ct^n$ on these fixed rescaled balls.
For a finite dyadic sum, the volume terms on the right telescope
exactly to $A_a(8Sr_N)-A_a(16S)$.  The noncollapsed Bishop inequality
gives $\lim_{t\downarrow0}A_a(t)\le\omega_n$ at every center,
where $\omega_n$ is the Euclidean unit-ball volume.  The fixed-radius
volume convergence in the preceding section gives
$A_a(16S)\ge\omega_n-o(1)$ uniformly in $a\in B_4(p)$.
Since $\sum_i\kappa r_i^2\le4\kappa/3$, this proves
\eqref{pi:poisson-sum}.  In particular this argument does not
assume that the neighboring centers are regular.

For completeness, the bound for the Dirichlet solution can be obtained
with uniform constants directly from
\eqref{pi:dirichlet-sobolev}.  Work in units $r_i=1$.  Let $\tau$
solve $-\Delta\tau=1$ with zero boundary values on $B_{8S}(a)$.
If $Y(k)=\m(\tau>k)$, testing with $(\tau-k)_+$ and applying
Sobolev and H\"older gives, for $l>k$,
\[
 Y(l)\le C(l-k)^{-2d}Y(k)^{2d-1}.
\]
Put $k_j=K(1-2^{-j})$.  Since $2d-1>1$, choose $K$ using only
$C,d,Y(0)$ so that this recursion implies
$Y(k_j)\le Y(0)2^{-jd/(d-1)}$ by induction.  Then
$\tau\le K$.  The harmonic function $f_i^a+a_i\tau$ has the
nonnegative boundary values $q^a$, so the weak maximum principle
gives $f_i^a\ge-a_i\tau$.  Together with $h_i^a\ge0$ this gives
$0\le h_i^a\le q^a+a_i\tau\le C$.  Rescaling proves the first
bound in \eqref{pi:poisson-boundedness}.

 The local Poisson gradient estimate gives $|\nabla f_i^a|\le C$
throughout $B_{4S}(a)$ in these units, since the values and the
constant right side are bounded.  More explicitly, for each
$q\in B_{4S}(a)$ apply the estimate on $B_{S/4}(q)$; its
eightfold enlargement $B_{2S}(q)$ is contained in $B_{6S}(a)$,
strictly inside the Dirichlet domain $B_{8S}(a)$.
The estimate
and its uniform local dependence follow from \cite{Kell13,Jiang14};
the local test-function interpretation is supplied by
Lemma~\ref{co:local-calculus} (see also \cite{CDNZPSW}).
Apply Lemma~\ref{pi:positive-source} to $h_i^a$ on this same
rescaled Dirichlet domain.  Its source mass is at most $P_i(a)$.
This proves both estimates in \eqref{pi:poisson-errors}, including
their displayed powers of $r_i$.

It remains to bound the full shifted measure.  Choose a nonnegative
test cutoff $\chi$, equal to one on $B_{2Sr_i}(a)$ and supported
in $B_{4Sr_i}(a)$, with $|\Delta\chi|\le Cr_i^{-2}$.
Since $h=q^a-f_i^a$ and $|\nabla q^a|\le8Sr_i$, we have
\[
 \int\left|\tfrac12|\nabla f_i^a|^2-q^a\right|\dd\m
 \le\int\bigl(|\nabla q^a||\nabla h|
                  +\tfrac12|\nabla h|^2\bigr)\dd\m
 \le Cr_i^{n+2}P_i(a).
\]
The integral can be restricted to the cutoff support or taken over
the Dirichlet domain.  Testing the definition of $M$ and using
$\Delta q^a=a_i\m-\mu_i^a$ now gives
\begin{align*}
 \int\chi\dd\mathcal B_\kappa(f_i^a,1)
 &=\frac12\int|\nabla f_i^a|^2\Delta\chi\dd\m
   +\kappa\int\chi|\nabla f_i^a|^2\dd\m
   -(2a_i-n)\int\chi\dd\m\\
 &=(n-a_i)\int\chi\dd\m-\int\chi\dd\mu_i^a
   +\kappa\int\chi|\nabla f_i^a|^2\dd\m
   +O(r_i^nP_i(a))\\
 &\le Cr_i^nP_i(a).
\end{align*}
The first two terms in the penultimate line are nonpositive, and
the curvature term is bounded by $C\kappa r_i^{n+2}$, already
included in $Cr_i^nP_i(a)$.  Positivity of
$\mathcal B_\kappa(f_i^a,1)$ proves
\eqref{pi:poisson-shifted-mass}.
\end{proof}

\begin{remark}[The units of the Poisson estimates]
\label{pi:poisson-units}
Under $d'=r_i^{-1}d$, $\m'=r_i^{-n}\m$ and
$\widehat f_i^a=r_i^{-2}f_i^a$, one has
$\Delta'\widehat f_i^a=a_i$ and $\kappa'=\kappa r_i^2$.
The measure $\mathcal B_{\kappa'}(\widehat f_i^a,1)$ is
$r_i^{-n}\mathcal B_\kappa(f_i^a,1)$.  Consequently all the
Poisson bounds above have size $CP_i(a)$ in their natural rescaled
units.  In particular, by \eqref{pi:shifted-triangle}, differences
of two replacements at the same scale are harmonic and have
Bochner mass at most $C(P_i(x)+P_i(z))$ on any of the common
interior balls.  This estimate controls the full measure, including
its possible singular part.
\end{remark}

\subsection{Radial differentiation of the coordinate span}

The Poisson estimates provide a summable scalar budget. We next
fit the radial derivatives by affine combinations of the harmonic
coordinates and distribute the errors over neighboring scales.

We first record the weak cancellation used to replace a radial
derivative by a harmonic function.

\begin{lemma}[Weak radial cancellation]
\label{pi:weak-cancellation}
Suppose that $f$ has constant Laplacian and $a$ is harmonic on a
ball on which both functions have bounded gradients.  Put
$g_a=\langle\nabla f,\nabla a\rangle$.  For interior Sobolev tests
$\phi\in W^{1,2}_0$ one has
\begin{equation}
 \int\langle\nabla g_a,\nabla\phi\rangle\dd\m
 =2\int(\Hess f-g)(\nabla a,\nabla\phi)\dd\m.
 \label{pi:weak-cancellation-equation}
\end{equation}
The identity is local, with the right side interpreted on a
slightly larger interior ball if a cutoff is needed.
\end{lemma}

\begin{proof}
Local test regularity and the Hessian product rule give
\[
 \nabla g_a=(\Hess f)(\nabla a,\cdot)
             +(\Hess a)(\nabla f,\cdot).
\]
For a test function $\phi$, express each of the two integrated
Hessian terms by differentiating
$\langle\nabla a,\nabla\phi\rangle$ or
$\langle\nabla f,\nabla\phi\rangle$ and integrating by parts.
The terms with $\Hess\phi$ cancel, leaving
\begin{align*}
 &\int\Hess a(\nabla f,\nabla\phi)\dd\m
  -\int\Hess f(\nabla a,\nabla\phi)\dd\m\\
 &\hspace{1cm}
 =-\int\Delta f\,\langle\nabla a,\nabla\phi\rangle\dd\m
   +\int\Delta a\,\langle\nabla f,\nabla\phi\rangle\dd\m
 =0.
\end{align*}
The constancy of $\Delta f$ and harmonicity of $a$ are both used
in the last equality.  This proves the identity with $\Hess f$
on the right.  Subtracting the metric tensor changes its right
side by $2\int\langle\nabla a,\nabla\phi\rangle=0$.
Both sides are continuous in the $W^{1,2}_0$ norm, by the bounded
gradient factors and the local $L^2$ Hessian bounds.  Density
therefore extends the identity to the stated tests.
\end{proof}

\begin{lemma}[Harmonic affine improvement]
\label{pi:affine-improvement}
There is an integer $m\ge3$, depending only on $n,S$, such that,
with $\theta=2^{-m}$, the following holds for all sufficiently
late spaces, all $i$ and all $a\in B_4(p)$.  If $w$ is harmonic
on $B_{Sr_i}(a)$, then there is an affine combination $F$ of the
components of $u$ with
\begin{equation}
 \|w-F\|_{2,B_{S\theta r_i}(a)}^{\mathrm{av}}
 \le\tfrac12\theta^{3/2}
       \|w\|_{2,B_{Sr_i}(a)}^{\mathrm{av}}.
 \label{pi:affine-improvement-equation}
\end{equation}
\end{lemma}

\begin{proof}
In a Euclidean ball, interior derivative estimates and Taylor's
formula at the center give
\[
 \|w-w(0)-\nabla w(0)\cdot X\|_{2,B_{S\theta}}^{\mathrm{av}}
 \le C_n\theta^2\|w\|_{2,B_S}^{\mathrm{av}},
 \qquad \theta<1/4.
\]
Choose $m$ so that $C_n\theta^2<\theta^{3/2}/4$.
If the asserted conclusion failed along increasingly late spaces,
rescale by $r_i$, normalize the nonzero norm of $w$ to one,
and apply interior harmonic compactness \cite{AH18}.  The limit
is Euclidean harmonic on $B_S$, with norm at most one there.
The normalized coordinate columns converge locally uniformly to
orthonormal affine Euclidean coordinates, by
Lemma~\ref{co:harmonic-normalization}.  An affine combination
of these columns therefore approximates the affine Taylor
polynomial of the limit uniformly on $B_{S\theta}$.
Strong convergence on this smaller ball contradicts the strict
Euclidean estimate.  This contradiction gives a single late-space
threshold for every center, scale, and harmonic $w$.  If the
normalizing norm is zero, the assertion is immediate.
\end{proof}

\begin{lemma}[Summable radial invariance]
\label{pi:radial-invariance}
There are functions $E_i(a)>0$, for $a\in B_4(p)$, such that
\begin{equation}
 E_i(a)\ge P_i(a),\qquad
 \sup_a\sum_{i=0}^\infty E_i(a)=o(1),\qquad
 \inf_a E_i(a)>0\quad\hbox{for each fixed }i.
 \label{pi:pinching-summability}
\end{equation}
Let $T=R_i(y)^{-1/2}$ with $d(a,y)\le C_{\mathrm{sep}}r_i$, and put
$U=T(u-u(a))/r_i$.  In the metric $r_i^{-1}d$, with
$\widehat f_i^a=f_i^a/r_i^2$, there is a matrix $A_i^a(T)$
satisfying $\|A_i^a(T)-\Id\|=o(1)$ and
\begin{equation}
 \left\|\langle\nabla\widehat f_i^a,\nabla U\rangle
       -A_i^a(T)(U-U(a))\right\|_{W^{1,2}(B_{S/2}(a))}
 \le C\sqrt{E_i(a)}.
 \label{pi:radial-invariance-equation}
\end{equation}
The ball and all derivatives in this display use the rescaled
metric.  Arbitrary additive constants may be used in $U$.
\end{lemma}

\begin{proof}
\emph{Harmonic correction and affine fits.}
We first work in the unrescaled units of the current member of the
sequence and fix a center $a$ and a covector $b$.  Define
\[
 g_b=\langle\nabla q^a,\nabla(b\cdot u)\rangle,
 \qquad
 g_{j,b}=\langle\nabla f_j^a,\nabla(b\cdot u)\rangle.
\]
The Poisson error and \eqref{pi:coordinate-input} imply
\begin{equation}
 \|g_b-g_{j,b}\|_{2,B_{Sr_j}(a)}^{\mathrm{av}}
 \le Cr_j\sigma_j^a(b)\sqrt{P_j(a)}.
 \label{pi:raw-radial-error}
\end{equation}
By Lemma~\ref{pi:weak-cancellation}, the equation for $g_{j,b}$
on this ball has divergence source
$2(\Hess f_j^a-g)(\nabla(b\cdot u),\cdot)$.
Its $L^2$ norm is at most
$Cr_j^{n/2}\sigma_j^a(b)\sqrt{P_j(a)}$.
Solve this divergence equation with zero boundary values to obtain
$\xi_{j,b}\in W^{1,2}_0(B_{Sr_j}(a))$.  Testing with
$\xi_{j,b}$ and then applying Poincar\'e gives
\begin{equation}
 \|\xi_{j,b}\|_{2,B_{Sr_j}(a)}^{\mathrm{av}}
 +r_j\|\nabla\xi_{j,b}\|_{2,B_{Sr_j}(a)}^{\mathrm{av}}
 \le Cr_j\sigma_j^a(b)\sqrt{P_j(a)}.
 \label{pi:radial-harmonic-correction}
\end{equation}
In particular $g_{j,b}-\xi_{j,b}$ is harmonic on that ball.

Let $F_{j,b}$ be the $L^2(B_{Sr_j}(a))$ orthogonal projection
of $g_b$ onto the finite-dimensional space of affine combinations
of $u$, and set
\[
 e_j(b)=r_j^{-1}
       \|g_b-F_{j,b}\|_{2,B_{Sr_j}(a)}^{\mathrm{av}}.
\]
The affine space is nondegenerate by harmonic normalization, so
these projections exist uniquely as functions and depend linearly
on $b$.  Apply Lemma~\ref{pi:affine-improvement} to
$g_{j,b}-\xi_{j,b}-F_{j,b}$.  The two errors in
\eqref{pi:raw-radial-error} and
\eqref{pi:radial-harmonic-correction}, restricted to the smaller
ball, acquire only the fixed volume factor depending on $m,S$.
Dividing by $r_{j+m}=\theta r_j$ gives
\begin{equation}
 e_{j+m}(b)\le\theta^{1/2}e_j(b)
                  +C\sigma_j^a(b)\sqrt{P_j(a)}.
 \label{pi:fit-recurrence}
\end{equation}
The factor $1/2$ in the harmonic improvement leaves more than the
needed margin in the coefficient of $e_j(b)$.

\smallskip\noindent
\emph{Qualitative convergence.}
We will also use the uniform estimate
\begin{equation}
 e_j(b)\le\varepsilon\sigma_j^a(b),\qquad\varepsilon=o(1).
 \label{pi:qualitative-radial-error}
\end{equation}
Here are details concerning the distance gradients in this assertion.
In units $r_j=1$, put $q=q^a/r_j^2$, normalize the coordinate
column at $a$, and write
$V=R_j(a)^{-1/2}(u-u(a))/r_j$.  On every fixed interior ball,
$V$ converges to orthonormal Euclidean coordinates and
$q^a/r_j^2-|V|^2/2\to0$ uniformly.  The norms of the gradients
of these two functions have the same limiting squared integrals:
for the distance function use $|\nabla(d_a^2/2)|=d_a$, and for
$|V|^2/2$ use the strong Gram convergence.  Their cross term has
the same limit.  Indeed, for a compactly supported Lipschitz cutoff
$\chi$,
\[
 \int\chi\langle\nabla(q-|V|^2/2),\nabla(|V|^2/2)\rangle\dd\m
 =-\int(q-|V|^2/2)
       \operatorname{div}(\chi\nabla(|V|^2/2))\dd\m=o(1),
\]
because $\Delta(|V|^2/2)=\sum_k|\nabla V_k|^2$ and the relevant
gradients are uniformly bounded.  Thus the two gradients are
close in local $L^2$.  Taking scalar products with $\nabla V$
and using strong Gram convergence shows
\[
 \langle\nabla q,\nabla V\rangle-V\longrightarrow0
 \quad\hbox{in local }L^2.
\]
This supplies an affine candidate in the defining infimum for
$e_j(b)$.  A contradiction subsequence over the center, scale,
and unit normalized covector makes the estimate uniform, proving
\eqref{pi:qualitative-radial-error}.

\smallskip\noindent
\emph{Summable squared errors.}
Choose a positive number $\eta=o(1)$ with $\eta\ge\varepsilon^2$;
for example, add the reciprocal of the sequence index to
$\varepsilon^2$.  Fix $\gamma=1/8$.  Iterating
\eqref{pi:fit-recurrence} in each of the finitely many residue
classes modulo $m$, and using \eqref{pi:coordinate-input}, yields
\begin{equation}
 e_i(b)\le C\sigma_i^a(b)\sqrt{\widetilde P_i(a)},\qquad
 \widetilde P_i(a)=\eta2^{-\gamma i}
          +\sum_{j=0}^i2^{-\gamma(i-j)}P_j(a).
 \label{pi:smoothed-fit-error}
\end{equation}
To see the exponent explicitly, the coefficient of a source at
$j<i$ before squaring is bounded by
$C_m2^{-(1/2-a_0)(i-j)}$.  The initial error has the same
decay times $\varepsilon$.  Weighted Cauchy--Schwarz bounds the
square of the source sum by an exponential convolution of $P_j$.
Since $a_0=1/16$ and $\gamma=1/8<1/2-a_0$, weakening that
convolution gives exactly \eqref{pi:smoothed-fit-error}.  This
calculation also covers the finitely many initial indices.

\smallskip\noindent
\emph{Removing the constant term.}
The constants in the affine fits must be removed with a quantitative
bound.  Put $c_j=F_{j,b}(a)$.  On $B_{Sr_{j+1}}(a)$,
the difference of consecutive fits has averaged $L^2$ norm at
most $Cr_j(e_j(b)+e_{j+1}(b))$.  Uniform affine nondegeneracy,
or harmonic interior evaluation on this ball, therefore gives
\[
 |c_j-c_{j+1}|\le Cr_j(e_j(b)+e_{j+1}(b)).
\]
Moreover $|c_j|\le Cr_j\sigma_j^a(b)$.  Indeed the raw radial
function has this averaged size, its projection has no larger
norm, and affine evaluation at the center is uniformly bounded.
The growth exponent $a_0<1$ implies $c_j\to0$ as $j\to\infty$.
Consequently, using \eqref{pi:smoothed-fit-error},
\begin{align}
 \frac{|c_i|}{r_i\sigma_i^a(b)}
 &\le C\sum_{l\ge i}2^{-(1-a_0)(l-i)}
       \bigl(\sqrt{\widetilde P_l(a)}
                +\sqrt{\widetilde P_{l+1}(a)}\bigr).
 \label{pi:inward-constant-sum}
\end{align}
Fix $\beta=1/16$ and define, with a sufficiently large fixed
constant $C_E$,
\begin{equation}
 E_i(a)=C_E\sum_{l=0}^\infty
                  2^{-\beta|i-l|}\widetilde P_l(a).
 \label{pi:definition-E}
\end{equation}
Weighted Cauchy--Schwarz in \eqref{pi:inward-constant-sum}
then gives $|c_i|\le Cr_i\sigma_i^a(b)\sqrt{E_i(a)}$.
The two geometric convolution kernels have finite sums, so
\[
 \sup_a\sum_iE_i(a)
 \le C\left(\eta+\sup_a\sum_iP_i(a)\right)=o(1).
\]
Also $E_i(a)\ge P_i(a)$ after increasing $C_E$, and the term
$l=0$ gives the center-independent positive floor
$E_i(a)\ge C_E\eta2^{-\beta i}$.

\smallskip\noindent
\emph{Gradient control and normalization.}
Combining the value estimates already proved gives
\[
 \|g_{i,b}-(F_{i,b}-c_i)\|_{2,B_{Sr_i}(a)}^{\mathrm{av}}
 \le Cr_i\sigma_i^a(b)\sqrt{E_i(a)}.
\]
The difference $g_{i,b}-F_{i,b}$ satisfies the divergence equation
of Lemma~\ref{pi:weak-cancellation}, because $F_{i,b}$ is harmonic.
Test it against its product with a squared Lipschitz cutoff,
equal to one on $B_{Sr_i/2}(a)$.  Cauchy--Schwarz and absorption
give the interior energy estimate
\[
 \|\nabla(g_{i,b}-F_{i,b})\|_{2,B_{Sr_i/2}(a)}^{\mathrm{av}}
 \le C\sigma_i^a(b)\sqrt{E_i(a)}.
\]
Here the value term is divided by $r_i$, and the divergence source
has the bound used in \eqref{pi:radial-harmonic-correction}.

For $T=R_i(a)^{-1/2}$, apply these scalar estimates to its rows,
rescale, and stack the resulting affine fits.  Their constants
at $a$ have been removed, so they give a matrix $A_i^a(T)$ and
\eqref{pi:radial-invariance-equation}.  The qualitative radial
convergence proved above, the Poisson error, and affine
nondegeneracy imply
$\|A_i^a(T)-\Id\|\le C(\varepsilon+\sqrt{P_i(a)}+
\sqrt{E_i(a)})=o(1)$ uniformly.  Finally, replacing
$R_i(a)^{-1/2}$ by $R_i(y)^{-1/2}$ conjugates this matrix by
$R_i(y)^{-1/2}R_i(a)^{1/2}$.  This matrix and its inverse have
uniformly bounded norms by fixed-distance normalization
comparability.  Thus both the quantitative bound and
$A_i^a(T)-\Id=o(1)$ persist.  Translating $U$ changes neither
side when its center value is subtracted.
\end{proof}

\subsection{Two centers and the sharp coordinate projection}

In Euclidean space, the difference $q^z-q^x$ of the squared-distance
potentials is affine, with gradient $x-z$.  We need a
quantitative version of both facts for the equal-Laplacian Poisson
replacements: their harmonic difference must be close to an affine
coordinate function, and its coefficient must be close to the
coordinate difference of the centers.

Fix a configuration \eqref{pi:center-configuration}.  In this
subsection take $r_i$ as the length unit and replace the measure
by $r_i^{-n}\m$.  Write
\[
 f^x=f_i^x/r_i^2,\qquad f^z=f_i^z/r_i^2,
 \qquad U=R_i(y)^{-1/2}(u-u(z))/r_i,
 \qquad s^2=E_i(x)+E_i(z).
\]
All balls and derivatives until the final scale conversion are in
these units.  We write $M$ for the Bochner measure with lower-bound
parameter $\kappa r_i^2$ in these units, and $g$ for the rescaled
metric tensor.  The equal-source construction gives
$\Delta f^x=\Delta f^z=a_i$.  Put
\begin{equation}
 v=f^z-f^x,\qquad
 L_a w=\langle\nabla f^a,\nabla w\rangle\ (a=x,z),\qquad
 h=\langle\nabla v,\nabla U\rangle.
 \label{pi:pair-notation}
\end{equation}
The domain nesting in \eqref{pi:radii} gives
\[
 B_{8R}(z)\subset B_{S/2}(x)\cap B_{S/2}(z),
 \qquad
 B_{2L+4}(y)\subset B_{R/2}(z).
\]
All these inclusions are strict with a fixed margin, since
$d(x,z)\le2C_{\mathrm{sep}}$ and $S/2=16R$.

The commutator below first identifies $h$ with $U(x)-U(z)$ up to
$O_{L^2}(s)$.  The sharp projection then gives an affine function
$c_0+b\cdot U$ whose gradient differs from that of $v$ by
$O_{L^2}(s)$ on the required interior ball.  Averaging against the
normalized Gram matrix recovers $b=U(x)-U(z)+O(s)$.  These two
conclusions transfer the $Cs^2$ Bochner bound for $v$ to the
coordinate function in the direction $U(x)-U(z)$.

\begin{lemma}[The two-center commutator]
\label{pi:commutator}
On $B_{8R}(z)$ one has
\begin{equation}
 \Delta v=0,\qquad
 |\nabla v|\le C,\qquad
 M(v)(B_{8R}(z))\le Cs^2,
 \label{pi:pair-poisson-bounds}
\end{equation}
and
\begin{equation}
 \|h-(U(x)-U(z))\|_{L^2(B_{8R}(z))}\le Cs.
 \label{pi:constant-direction}
\end{equation}
\end{lemma}

\begin{proof}
The first assertions follow from the common value of the Laplacian,
the Poisson gradient bounds, and
\eqref{pi:shifted-triangle}; see
Remark~\ref{pi:poisson-units}.  In particular
\[
 \|\Hess v\|_{L^2(B_{8R}(z))}
 +\|\Hess f^z-g\|_{L^2(B_{8R}(z))}\le Cs.
\]
Let $A_x,A_z$ be the matrices in
Lemma~\ref{pi:radial-invariance} for the common normalization
$R_i(y)^{-1/2}$.  Put
\[
 \delta=U(x)-U(z),\qquad B=A_z-A_x,\qquad c=A_x\delta.
\]
Since $U(z)=0$, subtraction of the two radial-invariance identities
gives
\begin{equation}
 h=BU+c+e,\qquad
 \|e\|_{W^{1,2}(B_{8R}(z))}\le Cs.
 \label{pi:commutator-affine}
\end{equation}
All coordinate gradients are bounded on this ball, and the values
of $U$, $\delta$, and $c$ are bounded by constants depending on
the fixed radii.  Also $A_x,A_z=\Id+o(1)$.

Set $L_vw=\langle\nabla v,\nabla w\rangle$.  Differentiate
the radial identity $L_zU=A_zU+O_{W^{1,2}}(s)$ along $\nabla v$.
Its error remains $O_{L^2}(s)$ because $|\nabla v|$ is bounded,
so
\[
 L_vL_zU=A_zh+O_{L^2}(s).
\]
The Hessian product rule separately gives the exact identity
\[
 L_vL_zU-L_zh
 =\Hess f^z(\nabla v,\nabla U)
       -\Hess v(\nabla f^z,\nabla U).
\]
The last displayed Hessian bounds and the bounded gradient factors
turn its right side into $h+O_{L^2}(s)$.  Thus
\begin{equation}
 L_zh=(A_z-\Id)h+O_{L^2}(s).
 \label{pi:commutator-equation}
\end{equation}
Every derivative here is legitimate in $L^2$ by local test
regularity and the gradient-product rule; no classical
differentiability of the metric is being used.

Insert \eqref{pi:commutator-affine} on both sides of
\eqref{pi:commutator-equation}, and use the radial identity once
more to differentiate $BU$.  This gives
\[
 \bigl(BA_z-(A_z-\Id)B\bigr)U-(A_z-\Id)c
 =O_{L^2}(s).
\]
Affine nondegeneracy, for example on $B_R(z)$, bounds its linear
and constant coefficients separately:
\begin{equation}
 \|BA_z-(A_z-\Id)B\|+|(A_z-\Id)c|\le Cs.
 \label{pi:commutator-coefficients}
\end{equation}
Writing $C_z=A_z-\Id$, the first operator is
$B+BC_z-C_zB$.  Consequently
\[
 (1-2\|C_z\|)\|B\|
 \le\|B+BC_z-C_zB\|\le Cs.
\]
For sufficiently late spaces, $\|C_z\|<1/4$, and hence
$\|B\|\le Cs$.  Next
\[
 (A_x-\Id)c=(A_z-\Id)c-Bc=O(s).
\]
Since $c=A_x\delta$ and $A_x^{-1}$ is bounded, this yields
\begin{equation}
 \delta-c=-A_x^{-1}(A_x-\Id)c=O(s).
 \label{pi:delta-from-c}
\end{equation}
Finally \eqref{pi:commutator-affine}, the bounded coordinate
values, and \eqref{pi:delta-from-c} prove
\eqref{pi:constant-direction} on the full ball $B_{8R}(z)$.
\end{proof}

The next lemma supplies the norm bound on a ball larger than the
projection domain.  This is the step which permits normalization
of a potentially very small projected residual.

\begin{lemma}[Radial transport of an unnormalized residual]
\label{pi:radial-transport}
Work in the rescaled units above.  Suppose that $w$ is a bounded
Sobolev function on $B_{4R}(z)$ with $\|w\|_\infty\le C_1$,
and that
\begin{equation}
 L_zw=w+\ell\cdot U+k+e_w,
 \qquad \|e_w\|_{L^2(B_{4R}(z))}\le C_1s.
 \label{pi:transport-hypothesis}
\end{equation}
Put $\alpha=\|w\|_{L^2(B_R(z))}$.  If
$|\ell|+|k|\le C_1(\alpha+s)$, then
\begin{equation}
 \|w\|_{L^2(B_{2R}(z))}\le C(\alpha+s),
 \label{pi:transport-conclusion}
\end{equation}
where $C$ is independent of $\alpha,s$.  The constants may depend
on $C_1$ and the already fixed radii and coordinate bounds.
\end{lemma}

\begin{proof}
The Sobolev product rule and $\Delta f^z=a_i$ give
\begin{equation}
 \operatorname{div}(w^2\nabla f^z)
 =(a_i+2)w^2+wF,\qquad
 F=2(\ell\cdot U+k+e_w),
 \label{pi:transport-divergence}
\end{equation}
with $\|F\|_{L^2(B_{3R}(z))}\le C(\alpha+s)$.
The identity is a weak identity: boundedness of $w$, its Sobolev
regularity, and the bounded gradient of $f^z$ justify the product
and its testing with compactly supported Lipschitz functions.

Here is a particular cutoff with the required sign.  For $t\ge0$
let $k_R(t)$ equal one for $t\le2R$, decrease linearly to zero
on $[2R,3R]$, and equal zero for $t\ge3R$.  Choose
$A_0=2(n+4)/(R-1)$ and put
\[
 \phi(t)=
 \begin{cases}
 1,&0\le t\le R,\\
 e^{-A_0(t-R)}k_R(t),&R<t<3R,\\
 0,&t\ge3R.
 \end{cases}
\]
This is a bounded nonnegative Lipschitz function.  It is bounded
below by $e^{-A_0R}>0$ on $[0,2R]$, and on $R<t<3R$ it
satisfies $-\phi'(t)\ge A_0\phi(t)$.  Since
$n\le a_i\le n+1$ and $R>1$, the choice of $A_0$ implies
\begin{equation}
 -t\phi'(t)-(a_i+2)\phi(t)
 \ge |\phi'(t)|+\phi(t)
 \quad\hbox{for almost every }R<t<3R.
 \label{pi:transport-cutoff-sign}
\end{equation}
The derivative and the cutoff are bounded by fixed constants.

Write $t=d_z$ and
\[
 e_f=\langle\nabla d_z,\nabla f^z\rangle-d_z,
 \qquad
 \mathcal A=\int_{R<d_z<3R}(|\phi'|+\phi)w^2\dd\m.
\]
The Poisson gradient error gives
$\|e_f\|_{L^2(B_{3R}(z))}\le Cs$, because
$\nabla(d_z^2/2)=d_z\nabla d_z$.  Test
\eqref{pi:transport-divergence} against $\phi(d_z)$.
Using \eqref{pi:transport-cutoff-sign}, and moving the inner-ball
term to the right, gives
\begin{equation}
 \mathcal A\le C\alpha^2
   +\left|\int\phi wF\dd\m\right|
   +\left|\int_{R<d_z<3R}\phi'e_f w^2\dd\m\right|.
 \label{pi:transport-energy}
\end{equation}
The decisive coefficient error has the bound
\begin{align}
 \left|\int_{R<d_z<3R}\phi'e_f w^2\dd\m\right|
 &\le\|w\|_\infty
       \left(\int|\phi'|w^2\dd\m\right)^{1/2}
       \left(\int|\phi'|e_f^2\dd\m\right)^{1/2}\notag\\
 &\le Cs\sqrt{\mathcal A}
 \le\tfrac14\mathcal A+Cs^2.
 \label{pi:transport-square-error}
\end{align}
All integrals in the middle line are over the annulus.
For the $F$ term, the inner ball contributes at most
$C\alpha(\alpha+s)$, while the annulus contributes at most
$C\sqrt{\mathcal A}(\alpha+s)$.  Another application of
Young's inequality therefore gives
\[
 \left|\int\phi wF\dd\m\right|
 \le\tfrac14\mathcal A+C(\alpha^2+s^2).
\]
Substitution in \eqref{pi:transport-energy} proves
$\mathcal A\le C(\alpha^2+s^2)$.  The positive lower bound for
$\phi$ up to radius $2R$ now gives
\eqref{pi:transport-conclusion}.

In particular, the fixed bound for $\|w\|_\infty$ in
\eqref{pi:transport-square-error} is used before $w$ is divided
by $\alpha$.  The conclusion is quadratic in $s$ at the level
of squared norms.
\end{proof}

\begin{lemma}[Sharp approximation by the harmonic coordinate span]
\label{pi:sharp-projection}
For all sufficiently late spaces, there is a vector $b\in\R^n$
with $|b|\le C$ such that
\begin{equation}
 \|\nabla(v-b\cdot U)\|_{L^2(B_{2L+4}(y))}\le Cs.
 \label{pi:sharp-projection-equation}
\end{equation}
The constants and the late-space threshold are uniform in the
configuration \eqref{pi:center-configuration} and the dyadic scale.
\end{lemma}

\begin{proof}
Subtract a constant from $v$ so that $v(z)=0$; this does not
change any gradient, Laplacian, or Bochner measure.  Its bounded
gradient gives a uniform bound for its values on $B_{8R}(z)$.
Let $c_0+b\cdot U$ be its $L^2(B_R(z))$ orthogonal projection
onto the affine coordinate span, and put
\[
 w=v-c_0-b\cdot U,\qquad
 \alpha=\|w\|_{L^2(B_R(z))}.
\]
Affine nondegeneracy and the bounds for $v,U$ give
$|c_0|+|b|\le C$.  The function $w$ is harmonic on $B_{8R}(z)$
and has uniformly bounded values and gradients there.

The gradient product rule and the Poisson estimates yield
\begin{equation}
 \nabla(L_zv-v)
 =(\Hess f^z-g)(\nabla v,\cdot)
    +\Hess v(\nabla f^z,\cdot),
 \qquad
 \|\nabla(L_zv-v)\|_{L^2(B_{8R}(z))}\le Cs.
 \label{pi:euler-gradient}
\end{equation}
Poincar\'e on $B_{4R}(z)$, using the gradient on $B_{8R}(z)$,
therefore gives a constant $k_v$ such that
$L_zv-v=k_v+O_{L^2(B_{4R}(z))}(s)$.
Subtracting the projection, and using
$L_zU=A_zU+O_{W^{1,2}}(s)$, gives
\begin{equation}
 L_zw=w+\ell\cdot U+k+e_w,
 \qquad
 \|e_w\|_{L^2(B_{4R}(z))}\le Cs,
 \label{pi:projected-euler-equation}
\end{equation}
where $\ell^T=b^T(\Id-A_z)$ and $k=k_v+c_0$.
We need a bound for these coefficients in terms of the small
residual norm.  Caccioppoli for the harmonic function $w$ gives
\[
 \|\nabla w\|_{L^2(B_{R/2}(z))}\le C\alpha.
\]
The bounded gradient of $f^z$ and
\eqref{pi:projected-euler-equation} now bound the $L^2$ norm
of $\ell\cdot U+k$ on this smaller ball by $C(\alpha+s)$.
Affine nondegeneracy gives
\begin{equation}
 |\ell|+|k|\le C(\alpha+s).
 \label{pi:forcing-coefficient-bound}
\end{equation}
Lemma~\ref{pi:radial-transport} consequently applies and yields
\begin{equation}
 \|w\|_{L^2(B_{2R}(z))}\le C(\alpha+s).
 \label{pi:projection-outer-bound}
\end{equation}

We claim that $\alpha\le Cs$ with a uniform constant for all
sufficiently late spaces.  If no such constant and threshold
existed, choose a sequence of increasingly late spaces and
configurations with $\alpha/s\to\infty$.  The positive floor
in \eqref{pi:pinching-summability} makes $s>0$, so
$\widehat w=w/\alpha$ is defined along this sequence.
The norm on $B_R(z)$ is one, and
\eqref{pi:projection-outer-bound} bounds its norm on $B_{2R}(z)$
uniformly.  Interior harmonic compactness therefore gives local
uniform and strong $W^{1,2}$ convergence on $B_{3R/2}(z)$ to
a Euclidean harmonic function $w_\infty$; see \cite{AH18}.
Use a rotation so that the normalized columns $U$, centered at
$z$, converge to the Euclidean coordinate column $X$.

The coefficients $\ell/\alpha,k/\alpha$ are bounded by
\eqref{pi:forcing-coefficient-bound}; pass to a convergent
subsequence.  The error $e_w/\alpha$ tends to zero in $L^2$.
The gradients of $f^z$ converge strongly to the radial vector
field $X$: the Poisson gradient error tends to zero, since
$P_i(z)\le s^2\to0$, and the squared-distance gradient
convergence was proved in
Lemma~\ref{pi:radial-invariance}.  The normalized harmonic
gradients have uniform interior bounds and converge strongly.
Thus scalar products of these gradients pass to the limit in
\eqref{pi:projected-euler-equation}, which becomes
\begin{equation}
 X\cdot\nabla w_\infty
 =w_\infty+\ell_\infty\cdot X+k_\infty
 \quad\hbox{on }B_{3R/2}(0).
 \label{pi:limiting-euler-equation}
\end{equation}
One may first pass this identity against interior tests and then
use the smoothness of the Euclidean harmonic limit.

Uniform convergence on a neighborhood of $\overline B_R(0)$
and convergence of the measures imply
\[
 \|w_\infty\|_{L^2(B_R(0))}=1,\qquad
 \int_{B_R(0)}w_\infty\dd X=0,\qquad
 \int_{B_R(0)}w_\infty X_j\dd X=0\quad(1\le j\le n).
\]
The boundary sphere has zero Euclidean measure, so restriction
to the projection ball causes no loss in these limits.
Expand $w_\infty$ in its harmonic Taylor series at zero.
The operator $X\cdot\nabla-1$ multiplies the homogeneous term
of degree $q$ by $q-1$.  Equation
\eqref{pi:limiting-euler-equation} therefore forces all terms of
degree at least two to vanish.  Analyticity makes $w_\infty$
affine on the ball.  Its affine orthogonality would then give
zero norm, a contradiction.  This proves the claimed uniform
bound $\alpha\le Cs$.

Finally Caccioppoli on the projection ball gives
$\|\nabla w\|_{L^2(B_{R/2}(z))}\le Cs$.
Since $B_{2L+4}(y)\subset B_{R/2}(z)$ and the constant $c_0$
does not affect gradients, this is
\eqref{pi:sharp-projection-equation}.
\end{proof}

\begin{proposition}[Summable directional pinching]
\label{pi:directional-pinching}
For every fixed $L,C_{\mathrm{sep}}>1$, the functions $E_i$ may be chosen
as in \eqref{pi:definition-E}, with the uniform summability and
positive floors in \eqref{pi:pinching-summability}, so that the
following holds for all sufficiently late spaces.  Define
\[
 D_{i,y}=R_i(y)^{-1}
       \bigl[r_i^{2-n}M(u)(B_{Lr_i}(y))\bigr]R_i(y)^{-1}.
\]
For every configuration \eqref{pi:center-configuration},
\begin{equation}
 D_{i,y}[u(x)-u(z)]
 \le Cr_i^2\bigl(E_i(x)+E_i(z)\bigr).
 \label{pi:directional-estimate}
\end{equation}
The same functions work simultaneously for all such centers and
all dyadic scales.  All enlargement constants are fixed before
choosing the late-space threshold.
\end{proposition}

\begin{proof}
Continue first in the rescaled units used in the preceding
subsection, and choose $b$ from
Lemma~\ref{pi:sharp-projection}.  The normalization at $y$ is
exactly
\[
 \fint_{B_1(y)}
   (\langle\nabla U_j,\nabla U_k\rangle)_{j,k}\dd\m=\Id.
\]
Consequently
\begin{align*}
 \fint_{B_1(y)}h\dd\m
 &=b+\fint_{B_1(y)}
       \langle\nabla U,\nabla(v-b\cdot U)\rangle\dd\m
 =b+O(s).
\end{align*}
Here the gradient bound for $U$ and
\eqref{pi:sharp-projection-equation} control the last error.
Equation~\eqref{pi:constant-direction} then gives
\begin{equation}
 b=U(x)-U(z)+O(s)=\delta+O(s).
 \label{pi:coefficient-recovery}
\end{equation}

The harmonic function $w=v-c_0-b\cdot U$ has gradient energy
at most $Cs^2$ on $B_{2L+4}(y)$.  Choose a nonnegative test
cutoff $\zeta$, equal to one on $B_L(y)$ and supported in
$B_{2L+2}(y)$, with bounded Laplacian.  Testing the definition
of the nonnegative measure $M(w)$ gives
\begin{align}
 M(w)(B_L(y))
 &\le\frac12\int|\nabla w|^2\Delta\zeta\dd\m
    +\kappa r_i^2\int\zeta|\nabla w|^2\dd\m\notag\\
 &\le Cs^2.
 \label{pi:residual-mass}
\end{align}
Together with \eqref{pi:pair-poisson-bounds}, positivity and
polarization imply $M(b\cdot U)(B_L(y))\le Cs^2$.
The normalized mass $M(U)(B_L(y))$ has bounded trace by
Lemma~\ref{co:bochner-small}.  Therefore
\eqref{pi:coefficient-recovery} implies
\[
 M((\delta-b)\cdot U)(B_L(y))
 \le|\delta-b|^2\tr M(U)(B_L(y))\le Cs^2.
\]
Another positive quadratic-form triangle inequality gives
\begin{equation}
 M(\delta\cdot U)(B_L(y))\le Cs^2.
 \label{pi:normalized-directional-mass}
\end{equation}
In particular the $O(s)$ coefficient error has contributed
$O(s^2)$ mass, including the possibly singular part of the
Bochner measure.

We finish by recording the scale conversion explicitly.  Return
to the unrescaled units, set $T=R_i(y)^{-1/2}$ and
$\delta u=u(x)-u(z)$.  Under $d'=d/r_i$, $\m'=r_i^{-n}\m$
and $U=Tu/r_i$ up to an additive constant, the matrix-valued
measure transforms as
\begin{equation}
 M'(U)=r_i^{2-n}T M(u)T^T.
 \label{pi:matrix-measure-scaling}
\end{equation}
Indeed the squared gradient matrix of $U$ in the primed metric
is $TGT^T$, the Laplacian gains a factor $r_i^2$, and the lower
curvature parameter becomes $\kappa r_i^2$.  Meanwhile
$\delta=U(x)-U(z)=T\delta u/r_i$.  Hence, on the corresponding
balls,
\begin{align*}
 M'(\delta\cdot U)(B_L'(y))
 &=r_i^{-n}\delta u^TR_i(y)^{-1}
       M(u)(B_{Lr_i}(y))R_i(y)^{-1}\delta u\\
 &=r_i^{-2}D_{i,y}[\delta u].
\end{align*}
Applying \eqref{pi:normalized-directional-mass} and
$s^2=E_i(x)+E_i(z)$ proves
\eqref{pi:directional-estimate}.
\end{proof}

\begin{remark}[Use with the heat-normalized matrices]
\label{pi:heat-interface}
The construction of $E_i$ and the radial argument used $R_i$ for
the scalar coordinate norms and did not use the heat matrices.
Lemma~\ref{co:heat-matrices}, specifically
\eqref{co:QR-comparison}, identifies $Q_i$ and $R_i$ up to
uniform relative $o(1)$.  They may therefore be interchanged in
bounded norm comparisons in the next section.  This introduces
only multiplicative constants; it adds no per-scale error term
to the summable quantities in
\eqref{pi:pinching-summability} or
\eqref{pi:directional-estimate}.  In particular the explicit
floor $E_i(a)\ge C_E\eta2^{-\beta i}$ is available when choosing
a single center that nearly minimizes $E_i$ on an active
coordinate patch.
\end{remark}
